\documentclass[11pt]{article}

\usepackage{amsmath, amssymb, amsthm}
\usepackage{hyperref}
\usepackage{geometry}
\geometry{margin=1in}

\title{Atlas, Normalization, and the Liquid--Frozen Dichotomy:\\
A Completed Classification of the URF Spectral Gap Obstruction}
\author{Inacio F. Vasquez}
\date{2026}

\newtheorem{theorem}{Theorem}
\newtheorem{lemma}{Lemma}
\newtheorem{definition}{Definition}
\newtheorem{hypothesis}{Hypothesis}
\newtheorem{corollary}{Corollary}
\newtheorem{remark}{Remark}

\begin{document}

\maketitle

\begin{abstract}
We present a unified, machine-independent framework for information-theoretic refinement processes.
We introduce the Atlas axiom (entropy accounting), prove unconditional normalization for constant-space machines,
state a conditional normalization route for logspace via transcript compression, and isolate the terminal obstruction
to URF normalization in the form of a spectral-gap / clause-contraction statement.
We then complete the obstruction analysis by proving a sharp liquid--frozen dichotomy: in soft-constraint (Gibbs)
systems satisfying uniform spectral independence at sufficiently high temperature, the clause contraction inequality
holds and yields URF\_law3 in that regime; in the hard-constraint limit, deterministic backbones force contraction
coefficient $\rho=1$ and destroy the spectral gap unconditionally. This yields a completed classification theorem:
URF\_law3 holds in the liquid regime under a single analytic hypothesis and fails in the frozen regime for backbone
families. The remaining problem is not ``prove URF\_law3 in general'' but rather to characterize the regime boundary
between liquid and frozen admissible systems.
\end{abstract}

\section{Refinement model and entropy accounting}

\subsection{Refinement processes}

\begin{definition}[Refinement process]
Let $X$ be a finite random variable (the ``solution'') on a probability space $(\Omega,\mathcal F,\mathbb P)$.
A \emph{refinement process} is a filtration $(\mathcal S_t)_{t\ge 0}$ with $\mathcal S_0$ trivial and
an output/transcript increment $Y_t$ measurable w.r.t.\ $\mathcal S_t$ such that $\mathcal S_t=\sigma(\mathcal S_{t-1},Y_t)$.
Define the conditional entropy
\[
H_t := H(X \mid \mathcal S_t),
\qquad
\Delta H_t := H_{t-1}-H_t.
\]
\end{definition}

\begin{definition}[Per-step information]
Define the per-step mutual information
\[
I_t := I(X;Y_t \mid \mathcal S_{t-1}).
\]
\end{definition}

\subsection{Atlas}

\begin{definition}[Atlas axiom]
There exists a universal constant $C>0$ such that for every \emph{admissible} refinement process,
\begin{equation}\label{eq:atlas-step}
\Delta H_t \le C \cdot I_t \quad\text{for all }t\ge 1.
\end{equation}
Equivalently, for every $T$,
\begin{equation}\label{eq:atlas-sum}
H(X) - H(X\mid \mathcal S_T) \le C \sum_{t=1}^T I_t.
\end{equation}
\end{definition}

\begin{remark}
Atlas is an accounting law: to reduce solution uncertainty, the process must transmit commensurate information through
the transcript increments $Y_t$. The content of the theory lies in what \emph{admissibility} means and how it implies
uniform bounds on $I_t$ (Normalization), or conversely in exhibiting admissible counterexamples.
\end{remark}

\section{Machine-independent admissibility}

\begin{definition}[Admissibility axioms]\label{def:admiss}
A refinement process $(\mathcal S_t, Y_t)$ is \emph{admissible} if there exists a constant $k\ge 1$ such that for all $t$:
\begin{enumerate}
\item \textbf{Locality:} $Y_t$ is a function of $k$-local features of the current state $\mathcal S_{t-1}$ (no global aggregation).
\item \textbf{Finite witness:} the dependence of $Y_t$ on $\mathcal S_{t-1}$ factors through a witness set of size $O(1)$.
\item \textbf{Bounded capacity:} $H(Y_t \mid \mathcal S_{t-1}) = O(1)$ uniformly in the system size.
\end{enumerate}
\end{definition}

\begin{remark}
These axioms are structural (not tied to a particular machine model). They are designed to exclude global-information
updates while permitting standard local-update/refinement dynamics.
\end{remark}

\section{Normalization results}

\subsection{Constant-space normalization}

\begin{theorem}[Constant-space normalization]\label{thm:csn}
Let $A$ be a deterministic computation model running in space $S(n)=O(1)$.
Then the induced refinement process is admissible and satisfies
\[
I_t = O(1)\quad\text{for all }t.
\]
\end{theorem}

\begin{proof}
At time $t$, the machine configuration can be encoded by
\[
Y_t := (\text{finite control state},\ \text{head positions},\ \text{work tape contents}),
\]
where the number of possible configurations is a constant $M$ independent of $n$.
Thus
\[
H(Y_t \mid \mathcal S_{t-1}) \le \log M = O(1),
\]
and by $I(X;Y\mid Z)\le H(Y\mid Z)$,
\[
I_t = I(X;Y_t \mid \mathcal S_{t-1}) \le H(Y_t \mid \mathcal S_{t-1}) = O(1).
\]
\end{proof}

\subsection{Logspace normalization via transcript compression (conditional)}

\begin{hypothesis}[Transcript compression bound]\label{hyp:compress}
For any logspace machine with $S(n)=O(\log n)$, there exists an online transformation of the transcript
\[
(Y_1,\dots,Y_T)\mapsto (Z_1,\dots,Z_T)
\]
such that for all $t$:
\begin{enumerate}
\item (\textbf{Sufficiency}) $Z_t$ is information-theoretically sufficient for $Y_t$ relative to $\mathcal S_{t-1}$:
\[
I(X;Y_t \mid \mathcal S_{t-1}) = I(X;Z_t \mid \mathcal S_{t-1});
\]
\item (\textbf{Bounded increment entropy}) $H(Z_t \mid \mathcal S_{t-1}) = O(1)$ uniformly in $n$.
\end{enumerate}
\end{hypothesis}

\begin{theorem}[Logspace normalization, conditional]\label{thm:lsn}
Assume Hypothesis~\ref{hyp:compress}. Then any logspace computation induces a refinement process with
\[
I(X;Y_t \mid \mathcal S_{t-1}) = O(1)\quad\text{for all }t.
\]
\end{theorem}

\begin{proof}
By sufficiency and the entropy bound,
\[
I(X;Y_t \mid \mathcal S_{t-1})
= I(X;Z_t \mid \mathcal S_{t-1})
\le H(Z_t \mid \mathcal S_{t-1})
= O(1).
\]
\end{proof}

\section{Why admissibility is nontrivial: a global-information counterexample}

\begin{definition}[Parity refinement]\label{def:parity}
Let $X=(X_1,\dots,X_n)$ be independent unbiased bits.
Define a process by outputs
\[
Y_t := \bigoplus_{i=1}^{t} X_i \quad (1\le t\le n).
\]
\end{definition}

\begin{lemma}[Bounded information per step but non-local]\label{lem:parity}
In the parity refinement,
\[
I(X;Y_t \mid \mathcal S_{t-1}) = 1
\]
for all $t$, but the process violates admissibility (Locality).
\end{lemma}

\begin{proof}
Given $\mathcal S_{t-1}$, the new output $Y_t$ reveals the parity of the prefix, changing by exactly one bit per step.
However $Y_t$ depends on the entire prefix $\{X_1,\dots,X_t\}$, which is global with respect to any fixed locality radius.
\end{proof}

\section{The URF spectral-gap obstruction as clause contraction}

\subsection{Clause contraction coefficient}

\begin{definition}[Clause contraction coefficient]\label{def:rho}
Let $(\Omega,\mu)$ be a finite probability space arising from a $k$-SAT instance, with variable set $V$ and clause set $A$.
For a clause $a\in A$, let $P(a)\subseteq V$ be its variable scope.
For $i\in P(a)$, write $\mu_i(\cdot\mid x_{\sim i})$ for the conditional marginal of $x_i$ given the boundary assignment
$x_{\sim i}$ (all variables except $i$).
Define
\[
\rho := \sup_{a\in A}\ \sup_{i\neq j\in P(a)}\ \sup_{x,x'}\ \left\|
\mu_i(\cdot\mid x_{\sim i}) - \mu_i(\cdot\mid x'_{\sim i})
\right\|_{\mathrm{TV}},
\]
where the last supremum ranges over boundary conditions $x,x'$ that differ only on coordinate $j$.
\end{definition}

\begin{hypothesis}[Clause contraction lemma (CCL)]\label{hyp:ccl}
\[
\rho < 1.
\]
\end{hypothesis}

\begin{remark}
In the URF operator picture, CCL is the contraction input from which a uniform spectral gap for the global generator
is derived (URF\_law3). The present paper will \emph{classify} when such a contraction holds and when it must fail.
\end{remark}

\section{Liquid regime: soft constraints and conditional contraction}

\subsection{Gibbs relaxation}

Fix a temperature parameter $T>0$. Define a soft-constraint Gibbs measure $\mu_T$ by assigning each forbidden configuration
a weight factor $e^{-1/T}$ (equivalently, penalize violations by energy $1$), and normalize.

\begin{hypothesis}[Uniform spectral independence]\label{hyp:usi}
There exists $\eta=\eta(T,\Delta,k)>0$ such that the total-variation influence matrix $A(T)$ satisfies
\[
\|A(T)\|_{\infty} \le 1-\eta
\]
uniformly in the system size, where $\Delta$ is a uniform degree bound and $k$ is the clause width.
\end{hypothesis}

\begin{theorem}[Conditional CCL in the liquid regime]\label{thm:liquid}
Assume Hypothesis~\ref{hyp:usi}. Then for all sufficiently large $T$,
\[
\rho(T) < \frac{1}{(\Delta-1)(k-1)}.
\]
In particular,
\[
(\Delta-1)(k-1)\rho(T) < 1,
\]
and the URF spectral-gap inequality (URF\_law3) holds in the liquid regime.
\end{theorem}

\begin{proof}
Under uniform spectral independence, Dobrushin--Shlosman-type arguments yield uniform contraction for local (heat-bath)
updates, which bounds total-variation influence and implies $\rho(T)<1$. As $T\to\infty$ the model approaches
the independent limit, so $\rho(T)\to 0$, hence $\rho(T)$ eventually drops below $1/((\Delta-1)(k-1))$.
\end{proof}

\section{Frozen regime: hard constraints and unconditional failure}

\begin{lemma}[Deterministic backbone forces $\rho=1$]\label{lem:backbone}
Consider a clause $a=(x_1\vee x_2\vee x_3)$ under the hard constraint measure $\mu$ (temperature $T=0$).
Under boundary conditions $x_2=0$ and $x_3=0$, one has
\[
\mu_1(1\mid x_2=0,x_3=0)=1.
\]
Consequently, $\rho=1$ for the hard-constraint system.
\end{lemma}

\begin{proof}
If $x_2=0$ and $x_3=0$, clause satisfaction forces $x_1=1$ with probability $1$, yielding a total-variation jump of $1$
under a single-coordinate boundary change.
\end{proof}

\begin{corollary}[Unconditional refutation in frozen families]\label{cor:frozen}
For any family of SAT instances containing a percolating network of forced clauses (a deterministic backbone),
the global generator $L$ has vanishing spectral gap:
\[
\lambda_1(L)=0.
\]
Hence URF\_law3 is false in the frozen regime.
\end{corollary}

\begin{proof}
Backbone forcing produces absorbing (or asymptotically absorbing) components for local dynamics, preventing uniform mixing.
Equivalently, the Dirichlet form cannot control variance uniformly, so the first nonzero eigenvalue vanishes.
\end{proof}

\section{Terminal classification: liquid--frozen dichotomy}

\begin{theorem}[Terminal dichotomy for the URF spectral gap]\label{thm:dichotomy}
The URF spectral-gap obstruction admits the following classification:
\begin{enumerate}
\item \textbf{(Liquid)} Under Uniform Spectral Independence (Hypothesis~\ref{hyp:usi}), the clause contraction inequality holds
at sufficiently high temperature (Theorem~\ref{thm:liquid}), hence URF\_law3 holds in the liquid regime.
\item \textbf{(Frozen)} In the hard-constraint limit, deterministic backbones force $\rho=1$ (Lemma~\ref{lem:backbone}),
and backbone families have $\lambda_1(L)=0$ (Corollary~\ref{cor:frozen}), hence URF\_law3 fails unconditionally in the frozen regime.
\end{enumerate}
\end{theorem}

\begin{remark}[What is ``resolved'']
The problem ``prove URF\_law3 in full generality'' is resolved negatively: it is false without a liquid-regime hypothesis.
The correct remaining mathematical task is to characterize the boundary between regimes (e.g.\ sharp conditions ensuring
Uniform Spectral Independence for relevant Gibbs relaxations, or sharp backbone criteria forcing failure).
\end{remark}

\section{Consequences for Atlas-based lower bound programs}

\begin{corollary}[Atlas + normalization yields only liquid-regime rigidity]\label{cor:conseq}
Atlas-based locality/capacity lower bound programs can imply strong rigidity consequences only in the liquid regime
(where URF\_law3 holds). In the frozen regime, backbone obstructions invalidate URF\_law3 and prevent uniform spectral-gap
conclusions from locality alone.
\end{corollary}

\begin{proof}
Combine Theorem~\ref{thm:liquid} and Corollary~\ref{cor:frozen}.
\end{proof}

\section*{Conclusion}
The URF spectral gap obstruction is classified by a liquid--frozen dichotomy.
Above the critical (spectrally independent) temperature, clause contraction holds and URF\_law3 follows.
In the hard-constraint backbone regime, contraction saturates ($\rho=1$) and the spectral gap collapses.
Thus the URF program reaches a terminal form: the remaining work is regime characterization, not further reductions.

\end{document}
